\section{Reproduction instructions}
\label{app:reproduction}

\subsection{Pinned versions}
\label{app:pinned_versions}

All results in this paper were produced against repository commit \texttt{db17ae5} on \texttt{master} (with the paper-only commits in \texttt{feat/geogen-paper} layered on top). The pilot dataset in \cref{sec:results} was scored by the verifier at that commit; we re-graded all 498 pilot records against the same commit's verification logic and confirmed zero \texttt{gate\_status} drift versus the as-collected JSONL records.

The renderer Docker image is pinned to the tag \texttt{tikz-renderer:db17ae5}, built from the \texttt{renderer/} directory at the same commit. The renderer was migrated mid-development from a DVI pipeline to a PDF+\texttt{fontspec} pipeline (master commit \texttt{a74970f}); reproductions before that commit will produce visually different SVGs but should yield identical strict-pass verdicts because the verifier operates on the SymPy-compiled coordinate set extracted from the TikZ source, not on rendered SVG bytes.

\subsection{Quickstart}

Clone the repository, install the Python deps with \texttt{uv sync}, start the renderer container:
\begin{lstlisting}
git clone https://github.com/<anon>/geogen
cd geogen
git checkout db17ae5             # paper-pinned commit
uv sync
docker run -d -p 8001:8001 \
    ghcr.io/<anon>/tikz-renderer:db17ae5
\end{lstlisting}

\subsection{Reproducing the headline numbers}

The headline pilot in \cref{sec:results} can be reproduced with:
\begin{lstlisting}
bash evals/leaderboard_launch.sh
\end{lstlisting}
This iterates over (model $\times$ strategy) combinations against \texttt{evals/scenarios\_pilot.yaml}, writing per-combo JSONL files into \texttt{evals/results/leaderboard\_pilot/}.
Subsequent analysis:
\begin{lstlisting}
uv run python -m evals.leaderboard_analyze \
    --input-dir evals/results/leaderboard_pilot_v3 \
    --output-md REPORT.md --output-csv leaderboard.csv

uv run python -m evals.leaderboard_plot \
    --input-dir evals/results/leaderboard_pilot_v3 \
    --output-dir docs/figures/geogen-pilot
\end{lstlisting}

\subsection{Re-rolling the dataset}

To generate a contamination-free re-roll of the templated split with a different seed:
\begin{lstlisting}
uv run python -m evals.generate_scenarios \
    --output evals/scenarios_generated_v2.yaml \
    --seed 42
\end{lstlisting}
The procedural variant generator preserves the 30-template / 200-per-tier structure and is the recommended defence against suspected contamination.

\subsection{Re-grading without LLM calls}

Because the pilot JSONL files store every model's TikZ output verbatim, the verifier can be re-run against a different commit's check semantics without spending another dollar on LLM tokens:
\begin{lstlisting}
uv run python -m evals.regrade \
    --input-dir evals/results/leaderboard_pilot_v2 \
    --output-dir evals/results/leaderboard_pilot_v3
\end{lstlisting}
We used this exact invocation to materialise \texttt{leaderboard\_pilot\_v3} from \texttt{leaderboard\_pilot\_v2} after pass~3 of \cref{sec:hardening} (the \propname{centroid} predicate): 10/498 records promoted, 0 demotions.
The same pattern audited verdict stability across the renderer/IR migration described in \cref{app:pinned_versions} (498 records, zero gate-status changes pre-pass-3), and is also the right tool for any future hardening pass: re-grade first, re-run only if the verdicts shift in a way the new code cannot reproduce.

\subsection{Docker images and dataset card}

We release:
\begin{itemize}
\item \texttt{ghcr.io/<anon>/tikz-renderer:db17ae5} --- the TikZ$\to$SVG renderer (LuaLaTeX + \texttt{fontspec} + \texttt{tkz-euclide}, pinned to the paper commit). A floating \texttt{:latest} tag tracks the current \texttt{master} for active development; reviewers should pull the pinned tag.
\item \texttt{ghcr.io/<anon>/geogen-eval:db17ae5} --- the full Python evaluator with all dependencies pre-installed.
\item \texttt{datasets/geogen-v1.0} on Hugging Face --- the 801-scenario benchmark in JSONL (\texttt{bench\_generated.jsonl} with 600 templated records and \texttt{bench\_curriculum.jsonl} with 201 curriculum records), generated deterministically from the YAML sources in \texttt{benchmark/definitions/} by \texttt{benchmark/definitions/build\_release\_jsonl.py}.
\end{itemize}

\subsection{Croissant metadata and datasheet}

A \textbf{Croissant 1.0} metadata file accompanies the release at \texttt{benchmark/definitions/croissant.json}.
It conforms to \texttt{http://mlcommons.org/croissant/1.0}, declares both released JSONL files with their SHA-256 checksums, and exposes per-record \texttt{id}, \texttt{prompt}, \texttt{tier}, and \texttt{tags} fields for both the \texttt{templated} and \texttt{curriculum} splits.
The file additionally carries the Croissant Responsible-AI vocabulary fields covering data collection, annotation, preprocessing, biases, limitations, recommended uses, and social impact (mined from the Gebru et al.\ datasheet, \texttt{docs/datasheet.md}, and the discussion sections of this paper).
Validation passes against the reference \texttt{mlcroissant} Python implementation v1.1.0:
\begin{lstlisting}
.venv/bin/python -c "
import mlcroissant as mlc
ds = mlc.Dataset(jsonld='benchmark/definitions/croissant.json')
print(ds.metadata.ctx.issues.report() or 'OK')"
\end{lstlisting}
The complete \GeoGen datasheet (\citealp{gebru2021datasheets} format) is at \texttt{docs/datasheet.md}; the Hugging-Face-ready dataset card is at \texttt{benchmark/definitions/README.md} and uses the same content.
